% Section 5 -- Results
% Prose drafted 28 October from the Option B outline (05-results-outline.tex,
% kept on disk, no longer typeset). Exhibits placed 28 October, each after the
% paragraph that first cites it: fig:lifecycle-fit, tab:payoff-horizon,
% fig:mechanisms-payoff (5.1; replaced Table tab:mechanisms-payoff on 30
% October, the table moved to Appendix F, A6-online-results.tex),
% tab:flows-accounting, tab:mechanisms-window
% (5.2; the two were one table, tab:mechanisms, until 29 October). Table bodies come from
% replication/output/{tables,figures}; captions and notes are draft (red).
\section{Results}\label{sec:results}

% Section~\ref{sec:data} left two explanations for why older owners live in
% less efficient homes. Efficiency may pay off over fewer years for an older
% owner, who is more likely to die before the cost of a renovation is
% recovered. We call this the payoff horizon. Or older households may buy
% homes less often, and a purchase is when most households get an efficient
% home. We call this the window. Section~\ref{sec:model} gave the direction
% in which each force moves the choices that create efficient homes, but not
% its size, so it could not say which one drives the age gradient. This
% section measures both at the estimates. The answer is that almost all of
% the gradient comes from purchases. The payoff horizon is short at every
% age, because efficiency wears off, and so it hardly differs between young
% and old owners. Purchases differ a great deal, and buying is how most
% efficient homes are acquired. Section~\ref{sec:results-horizon} measures
% the horizon and what it does, and Section~\ref{sec:results-window} the
% window. \open{Title: ``Results'', or one that states the question, e.g.\
% ``Why Older Owners Live in Less Efficient Homes''.}

 Section~\ref{sec:data} proposed two probable explanations for why older owners
live in less efficient homes. These were the  horizon channel and the purchase window. In Section~\ref{sec:model} the model was set up to include these two forces. This section asks
which of the two or whether their combination produces the age gradient at the estimates of
Section~\ref{sec:estimation}. We first show the lifecycle that the
estimates imply.

Figure~\ref{fig:lifecycle-fit} plots the lifecycle  profiles of homeownership,  the share of owners in an efficient home, mean consumption and mortgage status at the estimates.
Households rent when young and buy in their late twenties and thirties. Hence, the homeownership rate increases from 22 percent at 25 to 52 percent at 40 and
75 percent at 60 and remains constant at about 78 percent after that. Most of the ownership is through a mortgage. The share of households with a mortgage rises from about 20 percent to 60 percent by age 55, and decreases after that. By age 70, nearly three
quarters of households own outright. Consumption has the usual hump shape. It rises with income until the
late fifties and falls after retirement. The share of owners in an
efficient home peaks at 23 percent around 30, falls to 18 percent at 50,
and to 11 percent at 75. The share starts to fall soon after the ages at which
most households buy.


\begin{figure}[!h]
  \centering
  \caption[Lifecycle at the estimates]{{Lifecycle profiles at the estimates}}
  \label{fig:lifecycle-fit}
  \begin{subfigure}[t]{0.49\textwidth}
    \centering
    \includegraphics{fig-lifecycle-fit-own.pdf}
        \caption{{Homeownership}}
  \end{subfigure}\hfill
  \begin{subfigure}[t]{0.49\textwidth}
    \centering
    \includegraphics{fig-lifecycle-fit-abc.pdf}
        \caption{{Owners in an A--C home}}
  \end{subfigure}
  \par\medskip
  \begin{subfigure}[t]{0.49\textwidth}
    \centering
    \includegraphics{fig-lifecycle-fit-res.pdf}
    \caption{{Net income and consumption}}
  \end{subfigure}\hfill
  \begin{subfigure}[t]{0.49\textwidth}
    \centering
        \includegraphics{fig-lifecycle-fit-mort.pdf}
            \caption{{Owners by mortgage status}}
  \end{subfigure}
  \par\smallskip
  \begin{minipage}{0.96\textwidth}\footnotesize
    { \textit{Notes:} Simulated at the estimates of
    Table~\ref{tab:estimates}.
    Panel~(a) plots homeownership, overall and split by whether the home  is efficient (A--C) or not, as shares of all
    households. Panel~(b) plots  owners of an efficient home
    as a share of owners.
    Panel~(c) plots mean income and mean consumption.
    Panel~(d) plots owners with a mortgage and owners who own outright, as shares
    of all households.} 
  \end{minipage}
\end{figure}

% Road map (drafted 28 October, after the opening was rewritten). The draft
% group starts here, so the opening paragraph and Figure 5 print in black.
 Panel~(b) of Figure~\ref{fig:lifecycle-fit} shows how the efficient
share falls with age. However, from this we cannot deduce why this share is falling. As outlined above, there are two channels. The first
is the  horizon channel. Under this channel an efficient home pays off only while its owner lives
in it and it stays efficient. As survival probability reduces as households age, older owners have  fewer expected years left to enjoy the benefit of  an efficient home. The second  channel is the purchase window. A  buyer can choose an
efficient home for a small price premium whereas an owner who stays must pay
for the renovation. As the probability to buy reduces over age, this channel can lead to a declining efficient home share with age. In  the rest of
this section we measure the size of both these channels at estimated parameters. First, in  Section~\ref{sec:results-horizon}
we compute the payoff horizon by age and ask whether a larger or
longer-lasting payoff changes the gradient. Second, in Section~\ref{sec:results-window} we explore the purchase window. Here we follow owners into and out of efficient homes,
find the margin that produces the fall with age, and close the purchase window by
taking efficient homes off the market. Finally, we ask why older owners do not
renovate in place instead, and which factors determine how often they buy. We show that  almost all of the gradient comes from purchases. 


\subsection{\texorpdfstring{The horizon channel}{The horizon channel}}\label{sec:results-horizon}
  For the ease of exposition, we rewrite the expression of the payoff horizon $H_j$
\begin{align*}
  H_j =\sum_{k\geq 0} \beta^k \prod_{i=0}^{k-1} \phi_{j+i} \varrho\mu_{j+i+1}
\end{align*}
where $\beta$ is the discount factor, $\phi_{j+i}$ the probability that an
owner of age $j+i$ survives to the next year, $\varrho$ the probability that
an efficient home stays efficient, and $\mu_{j+i+1}$ the probability that an
owner of age $j+i+1$ does not move. Table~\ref{tab:payoff-horizon} computes the payoff horizon at different ages and illustrates the effect of these factors by adding one factor at a time. In column~(1) we keep discounting and survival only, setting $\varrho$ and
$\mu$ to one. In  column~(2) we add the loss of efficiency at the estimate
$\varrho=0.84$. Lastly in columns~(3) and~(4) we add moving, with $\mu_j$ either the model's
probability that an owner of an efficient home does not move or the
probability that an owner does not move in the Understanding Society survey.

With survival alone (column~(1)), the horizon falls steeply with age. At age 25, the horizon is 27 years which falls to ten years at age 75. This is the
intuition behind the first explanation. However, once we include the  loss of efficiency, the picture changes. At $\varrho=0.84$ efficiency lasts about six years on average, and
a home that is efficient today is still efficient five years later with
probability $\varrho^5\approx0.41$. Once this loss is added in column~(2), the horizon is
5.3 years at 25, 5.3 at 45, 5.0 at 65, and 4.4 at 75. In short, the horizon is essentially flat. Adding the model's probabilities of moving (column~(3)) lowers the horizon
to 3.4, 3.0, 3.1, and 2.9 years. With the moving probabilities of owners in
Understanding Society (column~(4)), the horizon is 3.8 years at 25, 4.7 at 45, 4.6 at 65, and 4.2 at 75.
As older owners move less, it offsets their higher probability of dying.\footnote{One caveat with these results is that $\varrho$ appears to be low and imprecisely estimated. The question arises: what if $\varrho$ were set to a higher value? In  Appendix~\ref{app:varrho-horizon} we recompute the horizon at
$\varrho=0.96$, at which efficiency lasts 25 years, with the other
parameters at the estimates. With loss of
efficiency alone, survival now matters and the horizon falls from 14.1 years at
25 to 7.8 at 75 (column~(5) of Table~\ref{tab:payoff-horizon-varrho}). Once we  include
moving it is again close to flat: 6.2 years at 25, 5.7 at 45, 6.1
at 65, and 5.4 at 75 with the model's probabilities of moving (column~(6)),
and 6.7, 9.8, 8.8, and 7.0 years with those in Understanding Society
(column~(7)). Hence, most owners move long before a rating that lasts 25 years wears
off. Section~\ref{sec:results-window} returns to this case once it has set
out the accounting.}

\begin{table}[t]
  \centering
  \caption[Payoff horizon of an efficient home, by age]{{Payoff
  horizon of an efficient home, by age}}\label{tab:payoff-horizon}
  \begin{threeparttable}
    \small
    \pkgtable{tab-payoff-horizon}
    \begin{tablenotes}[flushleft]\footnotesize
      \item {\textit{Notes:} The horizon $H_j$ of
      equation~\eqref{eq:payoff-horizon}, in years, for an owner of an
      efficient home at age $j$, at the estimates ($\beta=0.97$,
      $\varrho=0.84$). Each column adds one force: column~(1)
      has discounting and survival only, for an owner who never moves in a
      home that never loses efficiency; column~(2) adds the loss of A--C
      status at rate $1-\varrho$; columns~(3) and~(4) add moving, at the
      model's probability that an owner of an A--C home moves (3) or at the
      probability that an owner moves in Understanding Society (4). The last row is the horizon at
      75 as a percentage of the horizon at 25.}
    \end{tablenotes}
  \end{threeparttable}
\end{table}




The main message of Table~\ref{tab:payoff-horizon} is that the horizon is essentially flat once we take decay and moving into account. This has implications for choices for which horizon matters. In Section~\ref{sec:model-forces} we show that each
choice depends on a level of consumption. This is because the comfort of an efficient
home is worth more to a household that consumes more. So, an owner renovates in
place if its consumption is at least $c_j^*$
(equation~\eqref{eq:cstar-inplace}), and a buyer takes an efficient home if
its consumption is at least $c_b^*$ (equation~\eqref{eq:cstar-buy}). We call
the difference between the two the threshold gap. An owner whose consumption
lies in it will not renovate, but would buy an efficient home at its next
purchase. By equation~\eqref{eq:window}, the threshold gap is the extra
one time cost of renovating in place spread over the horizon.
If the horizon changes little with age, it implies that the threshold gap for a given home is about as wide
for an owner of 75 as for an owner of 35. Also, by
Proposition~\ref{prop:buy}, a higher probability of moving or dying raises
a buyer's threshold only in proportion to the transaction costs on the
premium. Therefore, the share of buyers who take an efficient home should not fall
much with age either. Hence, the two margins on which the horizon acts, what a
buyer chooses and whether an owner who stays renovates, should therefore be
close to flat in age. 



We next change the payoff of an efficient home by varying the comfort value $\eta$ and the probability of maintaining renovation status $\varrho$ and ask whether
the age profile of the efficient share shifts or tilts.
Figure~\ref{fig:mechanisms-payoff} plots the result. For each run we keep the other parameters at their
estimates and let households reoptimize. Each point in the figure plots the efficient
share of owners at 25--34 against the share at 65--75.  We superimpose dashed lines through the origin, along each of which the ratio of the two shares is constant. At the estimated point, the ratio is 1.75. A change that moves a point along its
line changes the level of the profile, not its slope. A change that moves it
toward the 45-degree line, on which young and old owners are alike, flattens
the gradient.

By varying  $\eta$ we set how large the payoff is each year by keeping its horizon fixed. If  young households choose the efficient home because  they expect to
enjoy it for a longer period, a larger payoff per year would be worth more to them. This would raise the share proportionally more among the young. However, it does not.
Without comfort the efficient share falls to 17.2 percent at 25--34 and 9.7
percent at 65--75. With comfort doubled it rises to 27.9 and 16.6 percent.
The resulting ratios are 1.76 and 1.68, against 1.75 at the estimates. Hence, the comfort
points stay close to the line of the estimates. In short, a larger payoff per year
actually raises the level of the profile, but does not alter its slope much.

Retention $\varrho$ sets how long the payoff lasts and it does so through two effects. First,
it lengthens the horizon and this should steepen the gradient, since survival
starts to matter once efficiency lasts.
Second, it slows the loss of efficiency in homes bought in the past, and this
can flatten the gradient. This is because older owners bought longer ago, so they lose
more of their efficiency when it wears off fast. The figure shows the net of the two
effects. Raising $\varrho$ from its estimate of 0.84 to 0.88 acts like a
larger payoff as  it lands almost exactly on the point of doubled comfort. However, raising it further,
to 0.92 and to 0.96, bends the path toward the  45-degree line, and the ratio
consistently falls from  1.75 at the estimates to 1.49 and 1.17.  The
gradient flattens, so the second effect wins. Purchases at 65--75 do not
change (Table~\ref{tab:mechanisms-payoff} in
Appendix~\ref{app:results-more}), so older owners do not catch up by buying
more often. When efficiency lasts, they catch up because they keep more of
the efficiency they bought.\footnote{Appendix~\ref{app:varrho-accounting} separates the two effects at
$\varrho=0.96$. The share of buyers who choose an efficient home now falls
with age, from 58 to 45 percent, which is the horizon at work. But older
owners lose less of the efficiency they bought when young, and this
outweighs it.}

 In conclusion, the  horizon channel does not produce the gradient. Once we take into account
efficiency wearing off and owners' moving probability, the horizon is about the same at every
age. Figure~\ref{fig:mechanisms-payoff} emphasizes this point. If young owners
chose efficient homes because they expect to enjoy them for longer, a larger
payoff each year would raise their share the most, and a payoff that lasts
longer would steepen the gradient. However, neither happens. A larger payoff raises
the share at every age proportionally. A payoff that lasts
longer actually flattens the gradient, as older owners then keep more of the
efficiency they purchased long ago. Therefore, what separates older owners from younger
ones is how long ago they bought as well as how often they buy.
Section~\ref{sec:results-window} takes this up.


\begin{figure}[t]
  \centering
  \caption[The efficient share of young and old owners when the payoff
  changes]{{The efficient share of young and old owners when the
  payoff changes}}\label{fig:mechanisms-payoff}
  \includegraphics[scale=0.7]{fig-mechanisms-payoff.pdf}
  \par\smallskip
  \begin{minipage}{0.96\textwidth}\footnotesize
    { \textit{Notes:} Each point is a simulation 
    that changes one parameter while keeping the others at the estimates of
    Table~\ref{tab:estimates} and allows for  households to reoptimize. A--C owners
    as a percentage of owners aged 25--34 (vertical) and 65--75
    (horizontal). Along each dashed line through the origin the share at 25--34
    is a fixed multiple of the share at 65--75, given at the end of the line;
    on the 45-degree line the two shares are equal. Comfort value
    $\eta$ takes three values: none, the estimate of 0.061, and twice the estimate. Retention
    $\varrho$ takes four values: the estimate of 0.84, 0.88, 0.92, and 0.96.}
  \end{minipage}
\end{figure}

\subsection{The purchase window}\label{sec:results-window}

 Section~\ref{sec:results-horizon} shows that the horizon is about the
same at every age, so the thresholds on which it acts barely change with
age. The choices themselves can still change, because whether a buyer takes
an efficient home, or an owner who stays renovates, also depends on how much
the household consumes, and consumption falls after retirement (panel~(c) of
Figure~\ref{fig:lifecycle-fit}). Table~\ref{tab:flows-accounting} therefore
measures the two choices directly. If both are flat in age, the gradient
must come from how often households buy, and the table checks this too. It follows
owners from one year to the next
and records how they enter and leave the stock of efficient homes, per 100
owners a year. Panel~A shows that at every age the main way in is to buy an
efficient home: 6.9 per 100 owners at 25--34, 4.1 at 35--64, and 1.9 at
65--75. Few buyers renovate at purchase, 0.3 to 0.4 per 100 owners.
Renovating in place adds 1.2 to 1.4 per 100 owners at every age. Buying is
therefore the only route that falls with age. As it falls, renovating in
place becomes a larger part of a smaller inflow: 14 percent of the inflow at
25--34, 24 percent at 35--64, and 36 percent at 65--75. This is why the
fixed cost of renovating, which a buyer of an efficient home avoids, weighs
most on older owners (Section~\ref{sec:identification}). On the way out,
the largest flow at every age is the loss of efficiency: 4.6, 3.5, and 2.3
per 100 owners.\footnote{ The loss is 18 to 21 percent of the stock in
the first row, more than $1-\varrho=16$ percent, because it applies to the
stock after the year's purchases and sales, which includes the homes bought
that year.} Selling to rent and selling to buy an inefficient home follow.
Death matters only at 65--75, and even there it removes 0.24 per 100
owners. The net inflow is 1.5 per 100 owners at 25--34, close to zero at
35--64, and $-0.5$ at 65--75. The stock of efficient homes grows among young
owners and shrinks among old ones, because purchases no longer replace what
wears off.

\begin{table}[!tp]
  \centering
  \caption[Flows into and out of efficient ownership, by age]{{Flows
  into and out of efficient ownership, by age}}\label{tab:flows-accounting}
  \begin{threeparttable}
    \small
    \pkgtable{tab-flows-accounting}
    \begin{tablenotes}[flushleft]\footnotesize
      \item {\textit{Notes:} Simulated at the estimates of
      Table~\ref{tab:estimates}. Panel~A: flows per 100
      owners in the age group a year. The stock is the A--C share of owners
      at the start of the year. Inflows are households outside the A--C stock
      at the start of the year. Lost A--C status applies to the stock after
      the year's purchases and sales, so as a share of the start-of-year
      stock it exceeds $1-\varrho$. Owners who move from one A--C home to
      another stay in the stock and are not shown. Panel~B: the inflow is
      $P(c_b+c_r)+N r$, where $P$ is the number of purchases by renters and
      owners of D--G homes per 100 owners, $c_b$ and $c_r$ are the shares of
      those purchases that are of A--C homes and that are renovated at
      purchase, $N$ is 100 less the stock, and $r$ is the in-place
      renovation rate of D--G owners. Panel~C: the change in each margin is
      weighted by the average of the other margins in the two age groups, so
      the parts add up to 100.}
    \end{tablenotes}
  \end{threeparttable}
\end{table}

To find which margin produces the fall in the inflow, we write it as
\begin{equation*}
  \text{inflow}=P\,(c_b+c_r)+N\,r,
\end{equation*}
where $P$ is the number of purchases by renters and owners of inefficient
homes per 100 owners, $c_b$ and $c_r$ are the fractions of those purchases
that are of efficient homes and that are renovated at purchase, $N$ is the
number of owners of inefficient homes per 100 owners, and $r$ is the
fraction of them that renovate in place. Purchases by owners of efficient
homes are left out, since they cannot add to the stock. Panel~B shows that
$P$ falls from 15.0 at 25--34 to 10.1 at 35--64 and 4.5 at 65--75.  The
two choices are close to flat, so the fall in consumption with age does not
undo the flat horizon. Between 40 and 46 percent of
purchases are of efficient homes at every age, and between 1.4 and 1.7
percent of owners of inefficient homes renovate in place each year. The
flat share of efficient purchases is the model's counterpart of Fact~3:
owners who bought recently live in similarly efficient homes at every age.
Panel~C decomposes the fall in the inflow from 25--34 to 65--75, 5.05 per
100 owners, and from 35--64 to 65--75, 2.52 per 100 owners. It evaluates the
change in each margin at the average of the other margins in the two age
groups, so the parts add up exactly to the fall. Purchases account for 99
percent of the first fall and 103 percent of the second. The other margins
offset one another, and together they account for 1 percent of the first
fall and $-3$ percent of the second. Owners over 65 renovate in place a
little less often than those aged 35--64, 1.4 against 1.7 percent, which
accounts for a tenth of the second fall. But older buyers more often
renovate at purchase, 6.4 against 4.2 percent of purchases, and with the
other margins this more than offsets it.\footnote{ The same conclusion holds for the stock. Rebuilding
the stock year by year from these margins, with purchases held at their rate
at 25--34, the efficient share at 65--75 would be 25.3 percent instead of
12.7. Holding any other margin at its rate at 25--34 changes it by at most
1.1 points. }

The accounting shows where the gradient appears, not why. Purchases are a
choice, and a short horizon is itself a reason not to move
(Section~\ref{sec:data}). We therefore change the choices themselves. The
most direct test closes the window. In the first experiment of
Table~\ref{tab:mechanisms-window}, efficient homes cannot be bought: a buyer can
still buy an inefficient home and renovate it at once, and an owner can
still renovate in place, and households reoptimize. The efficient share
falls to 6.9, 5.5, and 5.6 percent at 25--34, 35--64, and 65--75, and the gap
between the youngest and the oldest group shrinks from 9.6 points to 1.3.
The payoff horizon still differs by age in this economy, but without
efficient homes to buy almost no gradient remains. Almost all of the
gradient needs the window, which is the causal counterpart of the
accounting.

\begin{table}[t]
  \centering
  \caption[The efficient share of owners when the window changes]{{The efficient share of owners when the window changes}}\label{tab:mechanisms-window}
  \begin{threeparttable}
    \small
    \pkgtable{tab-mechanisms-window}
    \begin{tablenotes}[flushleft]\footnotesize
      \item {\textit{Notes:} Each row changes one element of the
      model, keeps the other parameters at the estimates of
      Table~\ref{tab:estimates} (values in parentheses), and lets households
      reoptimize. Shares are A--C owners as a percentage of owners in each
      age group, from 65,536 simulated households. The gap is the share at
      25--34 less the share at 65--75, in percentage points; the ratio is
      the share at 25--34 over the share at 65--75. No A--C homes for sale: a buyer
      cannot buy an efficient home, but can still buy an inefficient home and
      renovate it at once, and an owner can still renovate in place. EHS: the
      targets of Table~\ref{tab:estimates}.}
      \item {$^a$All home purchases, by renters and owners, per 100
      owners aged 65--75 a year.}
    \end{tablenotes}
  \end{threeparttable}
\end{table}

If owners who stop buying do not renovate in place instead, something must
make renovating in place costly relative to buying. In
equation~\eqref{eq:window} it is the extra one-off cost of renovating in
place, and the fixed cost $\xi_f$ is part of it. Halving $\xi_f$ narrows the
threshold gap: the efficient share rises at every age, most at 65--75 (15.0 against
12.8 percent), where more of the inflow comes from renovating in place, and
the ratio of the share at 25--34 to the share at 65--75 falls from 1.75 to
1.57. Without a fixed cost the ratio is 1.36. The fixed cost governs how
steeply the share falls with age, as Section~\ref{sec:identification}
argued, but even without it the threshold gap narrows and does not close.

What, then, sets how often older households buy? The last rows of
Table~\ref{tab:mechanisms-window} vary the strength of the bequest motive,
$\upsilon$, which is pinned down by the wealth of households over 65
(Section~\ref{sec:identification}), not by any moment of efficiency. With a
weak bequest motive, $\upsilon=5$ instead of the estimate
of 66.5, households aged 65--75 make 23.7 purchases a year per 100 owners
instead of 6.1, and the gradient reverses: the efficient share is 23.3,
28.6, and 30.8 percent at 25--34, 35--64, and 65--75. With a strong one,
$\upsilon=200$, they make 2.8 purchases per 100 owners, and the gap between
25--34 and 65--75 widens to 22.3 points. The share at 65--75 moves with
purchases: it rises by a factor of 2.4 when they rise by a factor of 3.9,
and about halves when they halve. This is what the window predicts. It
gives the bequest force of Section~\ref{sec:model-forces} a size, and it is
consistent with the bequest motive shaping the gradient through how often
the old buy. 

The accounting and the experiment without efficient homes for sale give the
same answer when efficiency lasts 25 years (Appendices~\ref{app:varrho-accounting} and~\ref{app:varrho-window}). At
$\varrho=0.96$, with the other parameters at the estimates, purchases account
for 90 percent of the fall in the inflow from 25--34 to 65--75 and 89 percent
of the fall from 35--64, against 99 and 103 percent at the estimates
(Table~\ref{tab:flows-accounting-varrho}). The horizon now has a role: the
share of purchases that are of efficient homes falls from 58 percent at
25--34 to 45 percent at 65--75, and this margin, on which the horizon acts,
accounts for 17 and 13 percent of the two falls. Without efficient homes for
sale, the gap between 25--34 and 65--75 shrinks from 6.2 points to 1.4,
against 9.6 to 1.3 at the estimates (Table~\ref{tab:mechanisms-varrho}). The
gap that remains without the window is about the same at both values of
$\varrho$. A longer-lasting payoff gives the horizon a role, but the window
still produces most of the gradient.

The gradient in the model is a purchase gradient. The payoff horizon is
short at every age, because efficiency wears off, so it hardly differs
between young and old owners, and the choices on which it acts hardly vary
with age. Older owners live in less efficient homes because they buy less
often, and buying is how most efficient homes are acquired. This is what
Fact~3 shows in the data. The distinction matters for policy. A subsidy that
lowers the cost of renovating in place narrows the threshold gap, while a
policy that acts at the point of sale works through the window. Section~\ref{sec:policy}
compares the two.

% Outline (kept for reference, not typeset):
% % Section 5 -- Results (outline: paragraphs and what each must say; no prose yet)
% Structure (28 October): an opening and two subsections, one per explanation
% of Section 2: the payoff horizon (5.1), then the window (5.2).
%
% Numbers from the draft exhibits of the replication package (R0, n = 65,536,
% CPU unless noted); none of the tables is in the paper yet:
%   lifecycle     output/figures/fig-lifecycle-fit.csv        (fig:lifecycle-fit)
%   horizon       output/tables/tab-payoff-horizon.csv        (tab:payoff-horizon)
%   flows and     output/tables/tab-flows-accounting.csv      (tab:flows-accounting;
%   accounting                                                 merged 28 October)
%   mechanisms    output/tables/tab-mechanisms.csv            (tab:mechanisms; GPU
%                 cells of the 30 September batch, decay rows on CPU)
%   varrho 0.96   data/model/r0/lifecycle/*-varrho0.96.csv    (flows, horizon)
% The search box is in data/model/r0/search-box.json.

\outpara{Opening: the question and the answer (before the first subsection)}
\begin{itemize}
  \item The question from Section~\ref{sec:data}: do older owners live in
    less efficient homes because efficiency pays off over fewer years
    (the payoff horizon), or because they buy homes less often (the
    window)? Section~\ref{sec:model} gave the direction of each force;
    this section gives their size at the estimates.
  \item The answer up front: almost all of the gradient comes from
    purchases. The payoff horizon is short at every age, because
    efficiency wears off, so it hardly varies with age. Older owners buy
    less often, and buying an efficient home is the main way to get one.
  \item Road map: the horizon (5.1), then the window (5.2).
  \item \open{Title: ``Results'' or one that states the question, e.g.
    ``Why Older Owners Live in Less Efficient Homes''.}
\end{itemize}

\outpara{Opening: the life cycle at the estimates (Figure fig:lifecycle-fit)}
\begin{itemize}
  \item The model's profiles of homeownership, the A--C share, mortgages,
    income, and consumption by age. Decided 12 and 21 October: model only,
    opens Section 5.
  \item What the reader needs from it: when households buy (renters buy in
    their late twenties and thirties, few purchases after 50), and the
    A--C share falling with age.
  \item Outcomes, not trade-offs: this is where sentences such as ``young
    households rent and save, then buy'' belong (kept out of
    Section~\ref{sec:model} on purpose).
\end{itemize}

\subsection{The payoff horizon}

\outpara{The horizon by age (Table tab:payoff-horizon)}
\begin{itemize}
  \item The horizon of equation~\eqref{eq:payoff-horizon}, built up one
    force at a time, at ages 25, 45, 65, 75.
  \item Survival alone: 26.9, 22.0, 14.6, 10.0 years. This is the intuition
    of the first explanation, and it falls steeply with age.
  \item Adding loss of efficiency: 5.3, 5.3, 5.0, 4.4. Adding moving at the
    model's rate: 3.4, 3.0, 3.1, 2.9. Loss of efficiency caps the horizon at
    every age, so survival hardly matters before 75.
  \item With the moving rates of owners in UKHLS the horizon rises in
    midlife (3.8, 4.7, 4.6, 4.2): the old move less, which offsets the
    higher probability of dying (Section~\ref{sec:model}).
\end{itemize}

\outpara{What a flat horizon implies}
\begin{itemize}
  \item Proposition~\ref{prop:window}: the gap between the two thresholds is
    the extra cost of renovating in place spread over $\widetilde H$. With
    $\widetilde H$ nearly flat in age, the gap is about as wide at 75 as at
    35. Proposition~\ref{prop:buy}: a buyer's choice of an efficient home
    hardly depends on age.
  \item So the margins on which the horizon acts (what a buyer chooses,
    whether an owner renovates in place) should be flat in age, and the
    horizon cannot produce much of a gradient. Section 5.2 checks this in
    the accounting.
  \item Optional: the reference-home example parked from Section 3.3 (bill
    saving against the cost of renovating) fits here, now that the
    estimates are known (parked-text.md).
\end{itemize}

\outpara{The payoff per year and how long it lasts (Table tab:mechanisms,
Panel A)}
\begin{itemize}
  \item Change one ingredient at a time, let households re-optimize, keep
    the other parameters at the estimates. Report the A--C share by age and
    the gap and the ratio between 25--34 and 65--75.
  \item Comfort (the payoff per year): none or doubled shifts every share
    (gap 7.4 and 11.3) but hardly the ratio (1.76 and 1.68). A larger
    payoff per year raises the level, not the age profile.
  \item Retention (how long the payoff lasts): $\varrho=0.88$, 0.92, 0.96
    raises every share; the gap widens and then narrows (11.4, 11.5, 6.2)
    and the ratio falls (1.68, 1.49, 1.17). A longer horizon does not
    steepen the gradient: with slow decay the stock remembers older
    purchases, so older owners catch up.
  \item \open{$\eta=0$ is the lower edge of the training box [0, 0.15],
    the same concern as for $\varrho$ at 0.80 and 0.99. Keep only the
    doubled cell (interior), add an interior cell such as $\eta/2$ (new
    run), or say so in the note.}
\end{itemize}

\outpara{When efficiency lasts 25 years: $\varrho=0.96$ (promised in
Section~\ref{sec:estimates})}
\begin{itemize}
  \item Why: $\varrho$ is imprecise and close to the lower edge of its
    range. At 0.96 efficiency lasts 25 years.
  \item The horizon roughly doubles (6.2 at 25, 5.4 at 75 at the model's
    moving rate; 6.7, 9.8, 7.0 with the moving rates in UKHLS) and survival
    starts to matter (with loss of efficiency alone, 14.1 at 25 against 7.8
    at 75).
  \item The fit is worse: shares of 43.1, 45.9, 36.8 against 21.0, 16.9,
    9.7 in the data.
  \item The test: do purchases still account for the fall in the inflow?
    \open{Not computed yet: the accounting at 0.96 needs only the existing
    simulation (flows\_by\_age, varrho 0.96), not a new run. The flow,
    horizon, and accounting tables at 0.96 go to the appendix.}
\end{itemize}

\subsection{The window}

\outpara{Where efficient homes come from (Table tab:flows-accounting,
Panel A)}
\begin{itemize}
  \item Per 100 owners a year, by age (25--34, 35--64, 65--75): the stock
    of A--C owners (22.5, 17.7, 12.7) and how it changes.
  \item Inflows by route. Buying an A--C home: 6.9, 4.1, 1.9. Renovating at
    purchase: 0.3, 0.4, 0.3. Renovating in place: 1.2, 1.4, 1.2. Buying is
    the main route at every age and the only one that falls with age.
  \item Because purchases fall, the share of the inflow that comes from
    renovating in place rises with age: 14, 24, and 36 percent. This is the
    claim Section~\ref{sec:identification} makes when it explains how the
    fixed cost shapes the profile.
  \item Outflows. Loss of efficiency is the largest outflow at every age
    (4.6, 3.5, 2.3), then selling to rent and selling to buy a D--G home;
    death matters only at 65--75 (0.24).
  \item Table note: loss of efficiency is 18--21 percent of the stock
    column, not $1-\varrho=16$ percent, because it applies to the stock
    after the year's purchases and sales (checked against the table).
\end{itemize}

\outpara{Which margin produces the fall (Panels B and C)}
\begin{itemize}
  \item Write the inflow as purchases per owner times the A--C share of
    purchases, plus the in-place renovation rate times the share of D--G
    owners. Decompose its fall from 25--34 and from 35--64 to 65--75.
  \item Purchases by renters and D--G owners, per 100 owners, fall from
    15.0 to 10.1 to 4.5 (A--C owners who sell and buy again are left out:
    they are already in the stock). They account for 99 percent of the
    fall from 25--34 and 103 percent from 35--64.
  \item The margins on which the horizon acts are flat, as 5.1 implies:
    the A--C share of purchases is 40--46 percent and the in-place rate
    1.4--1.7 percent of D--G owners at every age. The other margins offset
    one another: together they account for 1 percent of the fall from
    25--34 and $-3$ percent from 35--64 (Panel C). The flat A--C share of
    purchases is the model's version of Fact~3 (recent buyers are about as
    efficient at every age).
  \item With purchases held at the 25--34 rate, the A--C share at 65--75
    would be 25.3 percent instead of 12.7. Holding any other margin moves
    it by at most 1 point.
  \item Caveat: accounting is not causation. Purchases are a choice, and a
    short horizon is itself a reason not to move
    (Section~\ref{sec:data}). The next paragraph tests the window directly.
\end{itemize}

\outpara{Closing the window (Table tab:mechanisms, Panel B)}
\begin{itemize}
  \item With no efficient homes for sale, the gap falls from 9.6 to 1.3
    points and the shares to 6.9, 5.5, 5.6. The window is necessary for
    the gradient. This is the causal counterpart of the accounting.
\end{itemize}

\outpara{Why renovating in place does not replace buying}
\begin{itemize}
  \item The fixed cost is the extra cost of renovating in place in
    Proposition~\ref{prop:window}. Without it, or with it halved, every
    share rises and the profile flattens (ratio 1.36 and 1.57 against
    1.75): owners who no longer buy renovate in place instead. This is the
    role Section~\ref{sec:identification} gives $\xi_f$.
  \item \open{$\xi_f=0$ is the lower edge of the box [0, 0.2]; the halved
    cell is interior. Keep only the halved cell, or say so in the note.}
\end{itemize}

\outpara{What sets how often the old buy: the bequest motive}
\begin{itemize}
  \item A weak bequest motive ($\upsilon=5$): owners at 65--75 buy 23.7
    homes per 100 instead of 6.1, and the gradient reverses (gap $-7.5$).
    A strong one ($\upsilon=200$): they buy 2.8, and the gap widens to 22.3.
  \item The bequest motive matters through purchases, not through the
    value of the heirs' premium. This is the third force of
    Section~\ref{sec:model}, now with a size.
  \item \open{$\upsilon=5$ is the lower edge of the box [5, 500]. An
    interior weak value (e.g.\ 20) needs a new run.}
\end{itemize}

\outpara{Section close (one paragraph)}
\begin{itemize}
  \item Answer to Section~\ref{sec:data}: the gradient is a purchase
    gradient. Fact~3 is what this predicts.
  \item Bridge to Section~\ref{sec:policy}: a policy that lowers the cost
    of renovating in place acts on the margin that 5.2 shows cannot
    replace buying; one that acts at the point of sale acts on the window.
    Section~\ref{sec:policy} compares them.
\end{itemize}

\outpara{Exhibits (planning note, not text)}
\begin{itemize}
  \item Main text, four exhibits: Figure fig:lifecycle-fit (opening);
    Table tab:payoff-horizon (5.1); Table tab:flows-accounting (5.2; Panel
    A flows, Panel B margins, Panel C decomposition; merged 28 October,
    not yet in the paper); Table tab:mechanisms (Panel A comfort and $\varrho$,
    cited in 5.1; Panel B no efficient homes for sale, $\xi_f$, $\upsilon$,
    cited in 5.2; placed in 5.2, where most rows are used).
  \item Appendix: flows, horizon, and accounting at $\varrho=0.96$.
  \item \open{Before building tab:mechanisms: (i) one baseline: the cells
    were run on the GPU (22.4, 17.8, 12.8; gap 9.6), the flows on the CPU
    (22.5, 17.7, 12.7; gap 9.9); re-run the cells on the CPU or report both;
    (ii) ``purchases'' there are all purchases per 100 owners (6.0 at
    65--75), not the 4.5 of the accounting; use one definition; (iii) drop
    the cells at $\varrho=0.99$; the mandate cell goes to
    Section~\ref{sec:policy}.}
\end{itemize}

